\section{Introduction}
Marketplaces, banks and logistics firms keep records of the same businesses in several systems, and those
systems rarely share an identifier. One source lists ``Sharma Traders Pvt Ltd, Plot 14, MG Road, Pune'', a second
writes ``Sharma Traders Private Limited, 14 Mahatma Gandhi Rd'' and a third gives the name in Devanagari with
no PIN code. Deciding which records belong to the same real business is called entity resolution (ER) or
record linkage~\cite{fellegi1969,christen2012}. Mistakes are expensive in both directions: a missed link splits
one seller into two accounts, while a false link merges the payment history of two unrelated companies.

We study ER as a predictive modeling problem. For each record of a reference source (Source~1) we must predict the
set of records in two other sources that describe the same business. The data is the training set of the Amazon
ML Challenge 2026 business entity resolution task: 2.2 million Source~1 records from the US and India,
10.3 million Source~2/3 records and 7.6 million labelled matches. Our team took part in the challenge with a
two-stage gradient boosting system that reached a public leaderboard score of 0.959 (macro F0.5).
For this paper we return to the problem with the tools of a predictive analytics course and ask three questions.
How far do interpretable classifiers (logistic regression, linear and quadratic discriminant analysis) get
compared with gradient boosting on the same features? Which evidence drives a match, as measured by the estimated
logistic regression coefficients? And can a count model for the number of matches per entity (Poisson, negative
binomial, multinomial logit) support the decision of how many candidates to accept?

The contributions are: (i) a reproducible pipeline that turns the raw records into 43 pair features after
TF-IDF blocking; (ii) a comparison of four classifiers with grouped cross-validation, grid search and random
search, scored on a lockbox split that was never touched while building the challenge system; (iii) an
analysis of decision rules that convert pair probabilities into match sets, including an expected-F0.5 rule;
and (iv) a count-data view of the problem, which shows that match counts in this data are under-dispersed.
All code, the executed notebook, the result tables and a dashboard are in our GitHub repository\footnote{\url{https://github.com/Manthan-Vats/carcosa-entity-resolution}. The challenge data belongs to the organisers; the repository holds only a 1,500-entity demo sample for the dashboard.}.

\section{Related Work}
The probabilistic basis of record linkage goes back to Fellegi and Sunter~\cite{fellegi1969}, who scored record pairs
by likelihood ratios of field agreement and accepted pairs above a threshold. Surveys by Elmagarmid
et al.~\cite{elmagarmid2007} and Christen~\cite{christen2012} describe the now standard steps: normalisation,
blocking, pair comparison, classification and evaluation. Blocking limits the quadratic number of comparisons;
Papadakis et al.~\cite{papadakis2020} review the main families, of which token and q-gram based similarity search
is closest to our TF-IDF approach~\cite{salton1988}.

Pair comparison relies on string similarity. Cohen et al.~\cite{cohen2003} compared edit-based and token-based
metrics for names and found hybrid token measures strongest, while the Jaro-Winkler metric~\cite{winkler1990}
remains common for short names. Bilenko and Mooney~\cite{bilenko2003} showed that learning how to combine
similarities with a classifier beats hand-tuned rules. Magellan~\cite{konda2016} packaged this
features-plus-classifier recipe, where random forests and logistic regression are typical choices. Deep models
followed: DeepMatcher~\cite{mudgal2018} and Ditto~\cite{li2020ditto} embed both records with neural or
pre-trained language models and gain most on textual, dirty data. They need GPUs and their decisions are harder to
explain, so we keep to feature-based models, where gradient boosted trees~\cite{friedman2001,ke2017} are the strongest
standard learner and logistic regression~\cite{hosmer2013} offers direct inference on each piece of evidence.

The challenge metric is an F-measure averaged over entities, so the decision step matters as much as the
probabilities. Lipton et al.~\cite{lipton2014} showed that the F1-optimal threshold of a calibrated classifier is
half the optimal F1 score, and Jansche~\cite{jansche2005} trained models directly for expected F-measure.
We use the expected-F idea per entity rather than one global threshold.

\section{Data and Preparation}
\subsection{Records and labels}
Each record has an identifier, a business name, an address and a country label. Source~1 is deduplicated; every
Source~2/3 record matches at most one Source~1 entity, and a Source~1 entity can match zero, one or many records.
The noise is typical of commercial data: legal suffix variants (Pvt/Private, Corp/Corporation), word order changes,
typos, DBA names, landmark addresses (``near SBI ATM''), missing PIN codes and, for India, names written in
Devanagari. Table~\ref{tab:quality} summarises missing addresses and non-Latin names in our sample.

\subsection{Normalisation}
Names and addresses are lower-cased, ``\&'' becomes ``and'', and text in Indic scripts is mapped to Latin with a token
dictionary learned from matched training pairs (1,339 name tokens), backed by a generic transliteration library
(\texttt{anyascii}) for unseen tokens. House and plot numbers are extracted into a separate field, and their
leading zeros are stripped (``00604'' $\rightarrow$ ``604''). Missing values are not imputed: an empty address
stays empty and is exposed to the models as a binary feature, and number features take the value $-1$ when either
record has no number, together with an explicit missing-number indicator. No outliers were removed. Unusually long
names (the 99th percentile is about 50 characters) are real records, and the length features capture them.

\subsection{Study sample and splits}
Running the full pipeline on all records takes several hours, so the comparative study uses a 5\% sample: 5\%
of the Source~1 entities, all their true matches, and 5\% of the Source~2/3 records that match nobody. Because
each Source~2/3 record belongs to at most one entity, the sample keeps the ratio of matched to unmatched records
(4.7 pool records per entity in both). The sample has 110,341 entities and 625,153 records in total.

The split unit is the Source~1 entity, so all candidate pairs of an entity fall in the same split. We reuse the
fold assignment of our challenge system: five development folds f0--f4 (16\% of entities each), stratified by
country and number of matches, and a 20\% lockbox. Folds f1--f3 train the final models, f4 serves for early
stopping and threshold choice, f0--f4 together form the five groups of the cross-validation used for tuning, and
the lockbox is the test set.

\begin{figure}[t]
\centering
\begin{tikzpicture}[node distance=0.2cm, every node/.style={font=\footnotesize},
  box/.style={draw, rounded corners=2pt, align=center, minimum width=7.6cm, minimum height=0.5cm, inner sep=2pt, fill=gray!6},
  arr/.style={-{Stealth[length=2mm]}, thick}]
\node[box] (n) {\textbf{Normalisation}: lower-case, transliteration, house numbers};
\node[box, below=of n] (b) {\textbf{Blocking}: TF-IDF cosine top-$k$ (name, address, number views)};
\node[box, below=of b] (f) {\textbf{Pair features} (43): similarities, context, string metrics};
\node[box, below=of f] (c) {\textbf{Classifier}: LR / LDA / QDA / LightGBM $\rightarrow P(\text{match})$};
\node[box, below=of c] (d) {\textbf{Decision}: expected-F0.5 top-$k$ per entity};
\node[box, below=of d, fill=blue!5] (k) {\textbf{Count model}: Poisson / NB / multinomial logit for $|M(e)|$};
\draw[arr] (n) -- (b); \draw[arr] (b) -- (f); \draw[arr] (f) -- (c); \draw[arr] (c) -- (d);
\draw[arr, dashed] (k.east) -- ++(0.35,0) |- (d.east);
\end{tikzpicture}
\caption{Pipeline of the study. The count model is an alternative way to choose how many candidates to accept.}
\label{fig:pipeline}
\end{figure}

\section{Methodology}
\subsection{Blocking}
Comparing every Source~1 record with every pool record of the same country would mean about $10^{12}$ pairs on
the full data. Blocking keeps a short candidate list per entity. We build L2-normalised TF-IDF
vectors~\cite{salton1988} for four views: name character 3-grams (robust to typos), name words, address words and
house numbers. Very common terms are removed through a maximum document frequency (0.5\% for character
grams, 3\% for name words). For a weighted combination of views $v$ with weights $w_v$, the joint score of a
query $q$ and a pool record $r$ is
\begin{equation}
s(q,r)=\sum_{v} w_v \cos\!\big(\mathbf{x}^{(v)}_q,\mathbf{x}^{(v)}_r\big).
\end{equation}
Four searches with different weights (address-heavy, name-heavy, number-aware, and name-only among records without
an address) return their top 10--30 records, and the union forms the candidate set. Sparse top-$k$ matrix products
make this fast enough for a laptop CPU.

\subsection{Pair features}
Every candidate pair gets 43 features in four groups. \emph{Blocking similarities}: the four view cosines, the
joint score and the rank of the candidate in each search. \emph{Context}: the gap between a candidate's score and
the best score of the same entity, the gap to the best entity for the same candidate, ranks, and the number of
candidates per entity and per pool record. These features let a pointwise classifier see the competition among
candidates. \emph{String similarity}: Levenshtein ratio, token-sort and token-set ratios, partial ratio and
Jaro-Winkler on names~\cite{cohen2003,winkler1990}; token-set, token-sort and partial ratio on addresses; Jaccard
overlap and first-number agreement of house numbers. \emph{Record properties}: name and address lengths
(log-transformed), empty address, Devanagari script and candidate source.

For the linear and discriminant models we drop one feature of each pair with absolute Pearson correlation above 0.95
on the training folds, which removes redundant copies such as the joint score and its rank. Gradient boosting uses
all features.

\subsection{Classifiers}
\emph{Logistic regression} models the match probability of pair $i$ with feature vector $\mathbf{x}_i$ as
\begin{equation}
P(y_i=1\mid\mathbf{x}_i)=\sigma(\beta_0+\boldsymbol{\beta}^{\top}\mathbf{x}_i),\quad \sigma(z)=\frac{1}{1+e^{-z}},
\end{equation}
and the coefficients maximise the penalised log-likelihood
\begin{equation}
\ell(\boldsymbol\beta)=\sum_i\big[y_i\log p_i+(1-y_i)\log(1-p_i)\big]-\frac{1}{2C}\lVert\boldsymbol\beta\rVert_2^2 .
\end{equation}
For inference we refit an unpenalised model by maximum likelihood (Newton-Raphson) on standardised inputs, so that
$e^{\beta_j}$ is the odds ratio for a one standard deviation increase in feature $j$, and Wald $z$-tests give
standard errors and $p$-values~\cite{hosmer2013}.

\emph{Linear discriminant analysis} assumes Gaussian class densities with a shared covariance $\Sigma$ and classifies by
\begin{equation}
\delta_k(\mathbf{x})=\mathbf{x}^{\top}\Sigma^{-1}\boldsymbol\mu_k-\tfrac12\boldsymbol\mu_k^{\top}\Sigma^{-1}\boldsymbol\mu_k+\log\pi_k ,
\end{equation}
which gives a linear boundary~\cite{fisher1936,hastie2009}. We tune Ledoit-Wolf style shrinkage of
$\Sigma$~\cite{ledoit2004}. \emph{Quadratic discriminant analysis} allows a covariance $\Sigma_k$ per class, which
adds the term $-\tfrac12\log|\Sigma_k|$ and makes the boundary quadratic; its regularisation parameter shrinks
each $\Sigma_k$ towards the identity. \emph{LightGBM}~\cite{ke2017} fits an additive model of regression trees by
gradient boosting on the log-loss~\cite{friedman2001} and serves as the non-linear reference, since it is the model
of our challenge system.

\subsection{From probabilities to match sets}
The challenge score of entity $e$ with predicted set $\hat M$ and true set $M$ is
\begin{equation}
F_{0.5}(e)=\frac{(1+\beta^2)\,PR}{\beta^2P+R},\quad \beta=0.5,
\end{equation}
with precision $P=|\hat M\cap M|/|\hat M|$ and recall $R=|\hat M\cap M|/|M|$. An entity without matches scores 1
only when nothing is predicted, and the final score is the mean over entities. $\beta=0.5$ weighs precision twice
as much as recall. We compare three rules: a fixed threshold of 0.5; the threshold that maximises pair-level F1 on
f4; and an \emph{expected-F0.5 top-$k$} rule. With candidate probabilities sorted as $p_1\ge p_2\ge\dots$, the
rule picks
\begin{equation}
k^{*}=\arg\max_{k\ge 1}\frac{1.25\sum_{j\le k}p_j}{k+0.25\sum_j p_j},
\end{equation}
a ratio-of-expectations approximation of the expected F0.5, and predicts the empty set instead when
$\prod_j(1-p_j)$, the probability that the entity has no match among its candidates, is larger.

\subsection{Count models}
Let $N_e=|M(e)|$ be the number of true matches of entity $e$. From eight entity-level summaries of the candidate
list (counts of strong name and address agreements, the best joint and name scores, record lengths and country) we
fit a Poisson GLM with log link, $\log \mathbb{E}[N_e]=\mathbf{z}_e^{\top}\boldsymbol\gamma$, and a negative
binomial (NB2) model with $\operatorname{Var}(N_e)=\mu_e+\alpha\mu_e^2$, which relaxes the Poisson assumption
that the variance equals the mean~\cite{cameron2013}. A multinomial logit over the classes $\{0,1,2,3{+}\}$ treats
the count as a category. The predicted mean also gives a decision rule: accept the $\operatorname{round}(\hat\mu_e)$
highest-ranked candidates.

\subsection{Validation and tuning}
Hyperparameters are tuned by five-fold grouped cross-validation~\cite{kohavi1995} on 15,000 development entities
(about 790k pairs), with the predefined folds as groups so that no entity is split. Grid search covers the
logistic regression penalty $C\in\{10^{-3},\dots,10\}$ with and without class weighting, LDA shrinkage, and the QDA
regularisation parameter. For LightGBM, random search~\cite{bergstra2012} draws 12 settings of the number of leaves,
learning rate, minimum leaf size, row and column subsampling and L2 penalty. The tuning score is average precision,
because only about 6.5\% of the pairs are matches and accuracy or ROC-AUC alone would hide differences in
precision~\cite{saito2015}. Pair-level results are reported as AUC-ROC, PR-AUC, F1, precision, recall, log-loss and
Brier score; count models as MAE and RMSE. All models use scikit-learn~\cite{pedregosa2011},
statsmodels~\cite{seabold2010} and LightGBM.
